\title{xLSTM: Extended Long Short-Term Memory}

\begin{abstract}
The core mechanisms behind Long Short-Term Memory (LSTM) models, specifically the constant error carousel and gating structures, emerged in the 1990s. Over the years, LSTMs have proven to be robust and essential, notably forming the architecture for the earliest Large Language Models (LLMs). Nevertheless, with the introduction of the Transformer—characterized by highly parallel self-attention—a significant shift occurred as this new paradigm outperformed LSTMs at scale. This study seeks to answer: What results can be achieved in language modeling by scaling LSTMs up to the billion-parameter regime, employing modern techniques from current LLM research while addressing the inherent weaknesses of LSTMs? To this end, we propose exponential gating along with normalization and stabilization strategies. We also revise the conventional LSTM memory by developing two new variants: (i) sLSTM, which utilizes a scalar memory, scalar update, and incorporates a new memory mixing process, and (ii) mLSTM, introducing complete parallelism via matrix memory and a covariance-based update mechanism. By embedding these LSTM variants within a residual block structure, we construct xLSTM blocks which are subsequently layered to build comprehensive xLSTM architectures. The combination of exponential gating and the adapted memory design significantly enhances the performance of xLSTMs, enabling them to rival or exceed leading Transformers and State Space Models in terms of both efficiency and scalability.\\
Code available at: \url{https://github.com/NX-AI/xlstm}
\end{abstract}

\section{Introduction}
\label{sec:intro}

The Long Short-Term Memory (LSTM) ideas~\citep{Hochreiter:91,Hochreiter:97nips,Hochreiter:97}, i.e., the constant error carousel and gating, were introduced to overcome the vanishing gradient problem of recurrent neural networks~\citep{Hochreiter:91,Hochreiter:00book}:

\begin{align}
\label{eq:lstm_idea}
  c_t \ = \ f_t \ c_{t-1} \ + \ 
          i_t \  z_t \ , \quad  
          h_t \ = \ o_t \ \psi({c_t}) \ . 
\end{align}

The constant error carousel operates by adding cell inputs $z_t$ to the previous cell state $c_{t-1}$ (depicted in green), with this process regulated by sigmoid gates (shown in blue). Specifically, the input gate $i_t$ and the forget gate $f_t$ manage the modification of the cell state, whereas the output gate $o_t$ determines how the memory cell's contents are emitted, producing the hidden state $h_t$. Before being passed to the output gate, the cell state undergoes normalization or squashing through the function $\psi$, after which the gated output yields the hidden state.

Long Short-Term Memory networks (LSTMs) have achieved widespread adoption across a range of fields \citep{Hochreiter:01,Hochreiter:07,Schmidhuber:15}, dominating text generation tasks until the advent of Transformers in 2017~\citep{Vaswani:17}. Their utility has been substantiated in diverse sequential learning challenges, including text generation~\citep{Graves:13,Karpathy:15blog}, handwriting synthesis~\citep{Graves:13}, sequence-to-sequence translation~\citep{Sutskever:14nips}, program analysis~\citep{Zaremba:14arxiv}, automatic image captioning~\citep{Karpathy:15,Hossain:19}, source code generation~\citep{Karpathy:15blog}, and hydrologic modeling tasks such as rainfall-runoff prediction~\citep{Kratzert:18,Kratzert:19} and flood warning systems~\citep{Nearing:24}. In reinforcement learning, LSTMs consistently demonstrate superior performance as sequence models, powering solutions such as the AlphaStar agent in StarCraft II~\citep{Vinyals:17short}, OpenAI Five for Dota 2~\citep{OpenAI:19}, and controllers for nuclear fusion reactors~\citep{Degrave:22short}. LSTMs possess a remarkable ability for abstraction, allowing them to efficiently capture and maintain semantic features in their cell states~\citep{Karpathy:15blog}, as illustrated by the emergence of units dedicated to numbers and syntax~\citep{Lakretz:19}, linguistic structures~\citep{Bau:19}, and sentiment analysis~\citep{Radford:17arxiv}. Even today, LSTMs continue to underpin critical applications~\citep{Degrave:22short,Nearing:24} and remain a durable solution within the machine learning landscape.

Although LSTMs have achieved notable breakthroughs, they are hindered by three primary constraints: (i) Inability to update previously made storage selections. This shortcoming is illustrated using the {\it Nearest Neighbor Search} problem (see also Appendix~\ref{sec:appNearestNeighborSearch}): Given a reference vector, the task involves sequentially scanning a sequence to locate the most similar vector and subsequently retrieving its associated value at the end of the sequence. As depicted in the left panel of Figure~\ref{fig:lstmProblems}, the mean squared error is reported for this task. Standard LSTMs have difficulty replacing a stored entry when a more similar vector is encountered later in the sequence. By contrast, our proposed xLSTM addresses this issue using exponential gating. (ii) Restricted capacity for information storage—LSTMs require data to be compressed within scalar cell states. This is showcased with the {\it Rare Token Prediction} task. In the right panel of Figure~\ref{fig:lstmProblems}, we present token prediction perplexity from the Wikitext-103 corpus~\citep{Merity:17}, partitioned by token frequency. The LSTM’s constrained storage leads to inferior results on infrequent tokens, a deficit that our xLSTM overcomes via a matrix-based memory structure. (iii) Intrinsic lack of parallel computation, arising from memory mixing due to the recurrent hidden-to-hidden links between time steps, which necessitate strictly sequential operations.

Such drawbacks in LSTM architectures have driven the adoption of Transformers~\citep{Vaswani:17} for language modeling tasks. An important question arises: If we eliminate these constraints and scale LSTM-like architectures to match the dimensions of today’s Large Language Models, what level of language modeling performance can we attain?

\section{Extended Long Short-Term Memory}
\label{sec:xLSTM}

In order to address the constraints inherent to standard LSTM models, the Extended Long Short-Term Memory (xLSTM) framework introduces two principal innovations to the underlying LSTM paradigm described in Equation~\eqref{eq:lstm_idea}. The first is exponential gating, and the second is the development of novel memory structures, which together expand the traditional LSTM lineup with two additional variants: (i) sLSTM, discussed in Section~\ref{sec:sLSTM}, employs scalar memory representations, scalar updates, and a process of memory mixing; (ii) mLSTM, outlined in Section~\ref{sec:mLSTM}, operates with matrix-based memory and utilizes a covariance (outer product) update mechanism that is inherently suited for full parallelization. Both the sLSTM and mLSTM leverage exponential gating to improve upon conventional LSTM models. Notably, mLSTM forgoes the memory mixing (specifically, hidden-hidden recurrent links) in favor of parallel processing. Both sLSTM and mLSTM can be generalized to support multiple memory cells, with sLSTM specifically permitting memory mixing between its cells. Additionally, sLSTM may be structured with multiple heads, wherein memory mixing occurs only among cells within the same head and not across different heads, thus establishing a novel mode of memory mixing when combined with exponential gating. For mLSTM, the configuration of multiple heads is functionally indistinct from utilizing multiple cells.

When these advanced LSTM types are incorporated into residual block modules, they form what are termed xLSTM blocks, as detailed in Section~\ref{sec:xLSTMblock}. By stacking these blocks in a residual manner within deep learning architectures, one constructs xLSTM architectures, as described in Section~\ref{sec:xLSTMarchitecure}. Figure~\ref{fig:xlstm_sketch} presents a schematic depiction of the components comprising an xLSTM architecture.

\subsection{Review of the Long Short-Term Memory}
\label{sec:orgLSTM}

The foundational LSTM framework~\citep{Hochreiter:91,Hochreiter:97nips,Hochreiter:97} proposed a scalar memory cell functioning as a core unit for both computation and storage, thereby mitigating the issue of vanishing gradients~\citep{Hochreiter:91,Hochreiter:00book} via the constant error carousel enabled by the cell state transition mechanism. The architecture incorporates three distinct gates: the input gate, the output gate, and the forget gate, with the latter being added by \citet{Gers:00}. At each time step $t$, the LSTM memory cell is updated according to the following equations:

\begin{align}
c_t \ &= \  f_t \ c_{t-1} \ + \ i_t \ z_t &  & &\text{cell state} \\
h_t  \ &= \ o_t \ \tilde{h}_t \ , & \tilde{h}_t \ &= \ \psi \left( c_t\right)
  &\text{hidden state} \\
z_t \ &= \ \varphi \left( \tilde{z}_t \right) \ , 
  &\tilde{z}_t \ &=  \ \mathbf{w}^\top_{z} \ \mathbf{x}_t \ + \
  r_{z}  h_{t-1} \ + \  b_{z} \ \
  &\text{cell input} \\
i_t \ &= \ \sigma \left( \tilde{i}_t  \right) \ , 
  &\tilde{i}_t \ &= \ \mathbf{w}^\top_{i} \ \mathbf{x}_t \ + \
  r_{i}  \ h_{t-1} \ + \  b_{i} \ \
  &\text{input gate} \\
f_t \ &= \ \sigma \left(  \tilde{f}_t \right) \ , 
  &\tilde{f}_t \ &= \ \mathbf{w}^\top_{f} \ \mathbf{x}_t  \ + \
  r_{f}  \ h_{t-1} \ + \  b_{f} \ \
  &\text{forget gate} \\
o_t \ &= \ \sigma \left( \tilde{o}_t \right) \ , 
  &\tilde{o}_t  \ &= \ \mathbf{w}^\top_{o} \ \mathbf{x}_t \ + \
  r_{o}  \ h_{t-1} \ + \  b_{o} \ \
  &\text{output gate} 
\end{align}

The vectors $\mathbf{w}_{z}$, $\mathbf{w}_{i}$,
$\mathbf{w}_{f}$, and $\mathbf{w}_{o}$ serve as input weights, mapping the external inputs $\mathbf{x}_t$ to the cell input, input gate, forget gate, and output gate, respectively. The terms $r_{z}$, $r_{i}$, $r_{f}$, and $r_{o}$ denote the recurrent weights that link the previous hidden state $h_{t-1}$ with the respective cell input and gating mechanisms. Associated biases are given by $b_{z}$, $b_{i}$, $b_{f}$, and $b_{o}$. The functions $\varphi$ and $\psi$ describe the activations for the cell input and hidden state, typically using the $\tanh$ nonlinearity. In particular, $\psi$ functions to constrain or normalize the range of the cell state, preventing it from becoming unbounded. The gating units employ the sigmoid activation, represented as $\sigma \left( x \right)= 1/(1+\exp(-x))$.
Subsequent developments grouped several one-dimensional memory cells as a vector $\mathbf{c}_t \in \mathbb{R}^{d}$, enabling the implementation of recurrent weight matrices $\mathbf{R} \in \mathbb{R}^{d \times d}$, which facilitate mixing among the cell outputs~\citep{Greff:15} (see Appendix~\ref{sec:apporgLSTM} for elaboration).
Ablation experiments indicated the necessity of each individual element within the memory cell~\citep{Greff:15}.

\subsection{sLSTM}
\label{sec:sLSTM}

To enhance the capacity of LSTMs for adjusting storage choices, we integrate exponential gates (depicted in red) along with mechanisms for normalization and stabilization. Specifically, the input and forget gates are enabled to employ exponential activation functions. For normalization, a normalizer state is utilized, which accumulates the products of the input gate and subsequent forget gates over time.

The scalar sLSTM forward pass is:

\begin{align}
\label{eq:slstmforward}
c_t \ &= \  f_t \ c_{t-1} \ + \ i_t \ z_t & & 
 &\text{cell state} \\
n_t \ &= \  f_t \ n_{t-1} \ + \ i_t  & & 
 &\text{normalizer state} \\
h_t \ &= \ o_t \ \tilde{h}_t \ ,
  &\tilde{h}_t \ &= \ c_t / n_t
&\text{hidden state} \\
z_t \ &= \ \varphi \left( \tilde{z}_t \right) \ , 
  &\tilde{z}_t \ &=  \ \mathbf{w}^\top_{z} \ \mathbf{x}_t \ + \
  r_{z}  h_{t-1} \ + \  b_{z}
  &\text{cell input} \\
i_t \ &= \ \!\exp \left( \tilde{i}_t  \right) \ , 
  &\tilde{i}_t \ &= \ \mathbf{w}^\top_{i} \ \mathbf{x}_t \ + \
  r_{i}  \ h_{t-1} \ + \  b_{i}
  &\text{input gate} \\
f_t \ &= \ \sigma \left(  \tilde{f}_t \right) \ \text{OR} \
\!\exp \left(  \tilde{f}_t \right) \ ,
  &\tilde{f}_t \ &= \ \mathbf{w}^\top_{f} \ \mathbf{x}_t  \ + \
  r_{f}  \ h_{t-1} \ + \  b_{f}
  &\text{forget gate} \\
o_t \ &= \ \sigma \left( \tilde{o}_t \right) \ , 
  &\tilde{o}_t  \ &= \ \mathbf{w}^\top_{o} \ \mathbf{x}_t \ + \
  r_{o}  \ h_{t-1} \ + \  b_{o}
  &\text{output gate} 
\end{align}

We transfer the original LSTM gating techniques, i.e., input- and/or hidden-dependent gating plus bias term, to the new architectures. Exponential activation functions can lead to large values that cause overflows. Therefore, we stabilize gates with an additional state \mbox{$m_t$~\citep{Milakov:18arxiv}}:

\begin{align}
\label{eq:slstmstabil}
m_t \ &= \ \max \left( \log ( f_t ) + m_{t-1} , \log ( i_t ) \right) &\text{stabilizer state} \\
i'_t \ &= \ \!\exp \left( \log \left ( i_t \right) - m_t \right) \ = \ \!\exp \left( \tilde{i}_t - m_t \right) \ \, 
  &\text{stabil. input gate} \\
f'_t \ &= \ \!\exp \left( \log \left( f_t \right) + m_{t-1} - m_t \right) \ \, 
  &\text{stabil. forget gate}
\end{align}

In Appendix~\ref{sec:appsLSTM}, we demonstrate that substituting $f_t$ with $f'_t$ and $i_t$ with $i'_t$ during the forward computation leaves both the overall network output and the gradients of the loss relative to the parameters unaffected.

\paragraph{New Memory Mixing.} Similarly to standard LSTM architectures, sLSTM accommodates several memory cells (refer to Appendix~\ref{sec:appsLSTM}). The presence of multiple memory cells facilitates memory mixing via recurrent pathways, specifically $\mathbf{R}_{\mathbf{z}}$, $\mathbf{R}_{\mathbf{i}}$, $\mathbf{R}_{\mathbf{f}}$, and $\mathbf{R}_{\mathbf{o}}$ linking the hidden state vector $\mathbf{h}$ to the memory cell input $\mathbf{z}$ as well as the gates $\mathbf{i}$, $\mathbf{f}$, and $\mathbf{o}$. An innovative component in this memory mixing paradigm is the influence of exponential gating. In the updated sLSTM architecture, multiple heads are permitted, each executing memory mixing within its boundaries, but without inter-head interaction. The addition of heads, combined with exponential gating in sLSTM, provides a novel framework for memory mixing.

\subsection{mLSTM}
\label{sec:mLSTM}

In order to boost the storage capabilities of LSTM architectures, we generalize the traditional scalar memory cell $c \in \mathbb{R}$ to a matrix representation $\mathbf{C} \in \mathbb{R}^{d \times d}$. Accordingly, retrieval from memory proceeds through matrix multiplication. At each timestep $t$, we seek to record two vectors: a key $\mathbf{k}_t \in \mathbb{R}^d$ and its associated value $\mathbf{v}_t \in \mathbb{R}^d$, following the vocabulary of the Transformer model. Subsequently, at a future timestep $t + \tau$, the stored value $\mathbf{v}_t$ is recovered using a query $\mathbf{q}_{t+\tau} \in \mathbb{R}^d$. This approach corresponds to the framework of Bidirectional Associative Memories (BAMs) \citep{Kohonen:72,Anderson:72,Nakano:72,Anderson:77}.

The covariance update rule~\citep{Sejnowski:77,Dayan:91} for storing a key-value pair is

\begin{align}
\mathbf{C}_t \ &= \ \mathbf{C}_{t-1} \ + \ \mathbf{v}_{t} \ \mathbf{k}_t^\top \ .
\end{align}

We presuppose that a layer normalization precedes the mapping of inputs to keys and values, ensuring their mean is zero. The optimality of the covariance update rule~\citep{Dayan:91} lies in achieving maximal distinguishability among retrieved binary vectors, effectively maximizing the signal-to-noise ratio. Enhanced separability can be realized by restricting retrieval to pairwise associations and accepting quadratic computational costs, similar to attention mechanisms~\citep{Krotov:16,Krotov:17,Ramsauer:21}. Notably, the covariance update rule aligns with Fast Weight Programmers \citep{Schmidhuber:92ncfastweights,Schlag:21}, which subsequently incorporated a fixed decay rate applied to $\mathbf{C}_{t-1}$ and a fixed learning rate applied to 
$\mathbf{v}_{t} \mathbf{k}_t ^\top$~\citep{Ba:16}.
Building on this concept, we adapt the covariance update rule within the LSTM structure: here, the forget gate represents the decay rate, the input gate denotes the learning rate, and the output gate modulates the retrieved vector's magnitude.

For this matrix memory, the normalizer state is the weighted sum of key vectors, where each key vector is weighted by the input gate and all future forget gates. Again, the normalizer state keeps record of the strength of the gates. Since the dot product between query and normalizer state can be close to zero, we use the absolute value of this dot product and lower bound it by a threshold (typically 1.0) as done previously~\citep{Sun:23arxiv}. The mLSTM forward pass is:

\begin{align}
\label{eq:mlstm_recurrent_begin}
\mathbf{C}_t \ &= \  f_t \ \mathbf{C}_{t-1} \ + \ 
  i_t \ \mathbf{v}_t \ \mathbf{k}_t^\top & &  &\text{cell state} \\
\mathbf{n}_t \ &= \  f_t \ \mathbf{n}_{t-1} \ + \ 
  i_t \ \mathbf{k}_t & &  &\text{normalizer state} \\
\mathbf{h}_t  \ &= \ \mathbf{o}_t \ \odot \ \tilde{\mathbf{h}}_t
\ , \qquad \qquad 
\tilde{\mathbf{h}}_t \ = \   \mathbf{C}_t \mathbf{q}_t \ / \ 
  \max \left\{ |\mathbf{n}_t^\top \mathbf{q}_t|, 1 \right\} 
   &&&\text{hidden state} \\
\mathbf{q}_t \ &= \ \mathbf{W}_q \ \mathbf{x}_t \ + \ \mathbf{b}_q  & &  &\text{query input} \\
\mathbf{k}_t \ &= \ \frac{1}{\sqrt{d}} \mathbf{W}_k \ \mathbf{x}_t \ + \ \mathbf{b}_k  & &  &\text{key input} \\
\mathbf{v}_t \ &= \ \mathbf{W}_v \ \mathbf{x}_t \ + \ \mathbf{b}_v  & &  &\text{value input} \\
i_t \ &= \ \!\exp \!  \! \left( \tilde{i}_t  \right) \ , 
 \qquad \qquad \ \ \ \ \,
  \tilde{i}_t \ = \ \mathbf{w}^\top_{i} \ \mathbf{x}_t \ + \  b_{i} \ \ \ 
  &&&\text{input gate} \\
f_t \ &= \ \sigma \!  \left(  \tilde{f}_t \right) \ \text{OR} \ \!\exp \!  \! \left(  \tilde{f}_t \right) , \ \ \: \!
\tilde{f}_t \ = \ \mathbf{w}^\top_{f} \ \mathbf{x}_t  \ + \
  b_{f}
  &&& \text{forget gate} \\
\label{eq:mlstm_recurrent_end}
\mathbf{o}_t \ &= \ \sigma \left( \tilde{\mathbf{o}}_t \right) \ , \qquad \qquad \quad \ \ \ \,
  \tilde{\mathbf{o}}_t  \ = \ \mathbf{W}_{\mathbf{o}} \ \mathbf{x}_t \ + \
  \mathbf{b}_{\mathbf{o}}
  &&&\text{output gate} 
\end{align}

mLSTM architecture permits several memory cells, analogous to traditional LSTM. In mLSTM, having multiple heads or multiple cells yields comparable effects due to the absence of memory mixing. To maintain stability in mLSTM's exponential gating mechanism, we apply the stabilization methods utilized for sLSTM, as outlined in Equation~\ref{eq:slstmstabil}. 
Given that mLSTM does not mix memory across cells, its recurrence lends itself to a parallel implementation. Further information can be found in Appendix~\ref{sec:appmLSTM}.

\subsection{xLSTM Architecture}
\label{sec:xLSTMarch}

\paragraph{xLSTM Blocks.}
\label{sec:xLSTMblock}
An xLSTM block is designed to encode historical information in a high-dimensional, non-linear manner, thereby facilitating the separation of distinct histories or contextual backgrounds. Effectively distinguishing among previous states is essential for accurate sequence prediction, such as anticipating the next token. To achieve this, we leverage Cover’s Theorem~\citep{Cover:65}, which posits that patterns embedded non-linearly in a high-dimensional space are statistically more likely to be linearly distinguishable than those in their original, lower-dimensional form. We explore two types of residual block designs: (i) The post up-projection residual block, reminiscent of the Transformer architecture, where the historical context is non-linearly compressed in its original space before being projected linearly into a larger dimensional space, subjected to a non-linear activation, and subsequently mapped back linearly to the original dimension; this structure is illustrated in the left segment of Figure~\ref{fig:Backbones} and is shown in the third column of Figure~\ref{fig:xlstm_sketch}. An expanded depiction can be found in Figure~\ref{fig:appsLSTM_detailed} in the appendix. (ii) The pre up-projection residual block, analogous to State Space Model frameworks, initially applies a linear transformation to move into a higher-dimensional representation,

Instead of a linear approach, the summarization of prior information occurs through a non-linear transformation within a high-dimensional representation. This is subsequently followed by a linear projection back into the initial space. When an xLSTM block incorporates an sLSTM, we implement the post up-projection module. In contrast, for xLSTM blocks containing an mLSTM, the pre up-projection module is employed, reflecting the expansion in memory capacity found in higher dimensions. For further clarification, readers can consult the left-hand side of Figure~\ref{fig:Backbones}, the third column of Figure~\ref{fig:xlstm_sketch}, or alternatively, Figure~\ref{fig:appsLSTM_detailed} within the appendix.

\paragraph{xLSTM Architecture. }
\label{sec:xLSTMarchitecure}
The xLSTM network structure arises from the residual composition of foundational units~\citep{Srivastava:15,He:16}. Our approach utilizes the widely-adopted pre-LayerNorm~\citep{Ba:16b} residual architectures, which serve as the foundation for recent Large Language Models. Reference the rightmost column of Figure~\ref{fig:xlstm_sketch} for visual details.

\subsection{Memory and Speed Considerations}

In contrast to Transformer architectures, xLSTM networks exhibit linear computational cost and maintain constant memory requirements as sequence length increases. This compressive memory property of xLSTM makes it particularly advantageous for deployment in industrial environments and on edge devices.

Although the memory mechanism in mLSTM does not necessitate learnable parameters, it introduces a computational burden due to its reliance on a $d \times d$ matrix memory structure and corresponding $d \times d$ update process. This framework embodies a trade-off between memory capacity and computational demand. Nonetheless, since these computations are highly parallelizable on GPUs, they impart only a negligible impact on total processing time.

Similar to approaches such as FlashAttention~\citep{Dao:22,Dao:23} and GLA~\citep{Yang:23arxiv}, mLSTM operations can be efficiently parallelized. In contrast, the sLSTM model lacks this parallelizability because of hidden-to-hidden memory mixing. To address this, we engineered a highly optimized CUDA implementation, with GPU register-level enhancements, resulting in sLSTM running at no more than twice the computation time of mLSTM in typical scenarios.

\section{Related Work}
\label{sec:related}

\paragraph{Linear Attention.} Multiple approaches have been proposed to address the quadratic scaling of Transformer attention with respect to context length, aiming to reduce it to linear complexity. The Synthesizer model generates synthetic attention weights independently of interactions between tokens~\citep{Tay:20arxiv}. Linformer employs a low-rank projection to represent self-attention, effectively enabling linear approximation~\citep{Wang:20arxiv}. Linear Transformer transforms the attention mechanism into a form that allows linearization \citep{Katharopoulos:20}. Performer achieves a linear approximation of the attention softmax using a positive orthogonal random feature method~\citep{Choromanski:21}. Alternatives to attention, such as fast long convolutions, have been integrated into Structured Global Convolution (SGConv)~\citep{Li:22} and the Hyena Hierarchy~\citep{Poli:23}.

\paragraph{State Space Models.} State Space Models (SSMs) have gained significant traction in recent research due to their linear complexity with respect to context length and their competitive results vis-à-vis Transformers. Early developments in this area include the Structured State Space sequence model (S4)~\citep{Gu:21}. Subsequent models that expanded on this framework are the Diagonal State Space (DSS) model~\citep{Gupta:22}, Gated State Space (GSS) models~\citep{Mehta:22}, S5 model~\citep{Smith:22}, Bidirectional Gated SSM (BiGS)~\citep{Wang:22}, H3 model~\citep{Fu:23}, and Mamba~\citep{Gu:24arxiv}.

\paragraph{Recurrent Neural Networks.} Recurrent Neural Networks (RNNs) have been proposed as alternatives to Transformer and attention mechanisms, primarily because their computational cost scales linearly with sequence length. RNN architectures employing Deep Linear Recurrent Units (LRUs) have demonstrated strong performance in language modeling tasks~\citep{Orvieto:23, De:24arxiv}. Hierarchically Gated Linear RNN (HGRN)~\citep{Qin:23} and its successor HGRN2~\citep{Qin:24arxiv} have also yielded favorable outcomes in similar settings. Among RNN-based strategies for scaling language models, RWKV~\citep{Peng:23arxivshort,Peng:24arxivshort} stands out, as it offers results comparable to Transformer models.

\paragraph{Gating.}
The concept of gating, central to the architecture of LSTM, has been revisited and reframed in numerous recent methodologies. Variations and extensions of gating mechanisms appear in models such as HGRN~\citep{Qin:23}, HGRN2~\citep{Qin:24arxiv}, Gated Linear Attention (GLA)~\citep{Yang:23arxiv}, Gated State Space (GSS) models~\citep{Mehta:22}, Bidirectional Gated SSM (BiGS)~\citep{Wang:22}, Moving Average Equipped Gated Attention (MEGA)~\citep{Ma:22}, RWKV~\citep{Peng:23arxivshort}, and Mamba~\citep{Gu:24arxiv}.

\paragraph{Covariance Update Rule.}
To improve the memory capability of the mLSTM cell, a matrix-based memory system employing a covariance update rule was incorporated. Models implementing similar update procedures include Fast Weight Programmers~\citep{Schmidhuber:92ncfastweights,Schlag:21}, RWKV-5, RWKV-6~\citep{Peng:24arxivshort}, Retention~\citep{Sun:23arxiv}, Linear Transformer~\citep{Katharopoulos:20}, and HGRN2~\citep{Qin:24arxiv}.

\paragraph{Most Related.} The models most similar in principle to xLSTM include Retention~\citep{Sun:23arxiv}, RWKV~\citep{Peng:23arxivshort,Peng:24arxivshort}, and HGRN2~\citep{Qin:24arxiv}. These architectures utilize mechanisms such as matrix memory and gating. Unlike the newly introduced sLSTM, however, they lack support for memory mixing. This capability is crucial for addressing state tracking tasks, which positions LSTMs as more expressive compared to both State Space Models (SSMs) and Transformers~\citep{Merrill:24,Deletang:23}. State tracking is essential for tasks like code evaluation or maintaining context regarding entities throughout extended narratives.

\paragraph{Residually Stacking Architectures.} The xLSTM architectures, consistent with the majority of modern deep learning frameworks, are assembled by residual stacking of individual modules~\citep{Srivastava:15,He:16}.

This design paradigm facilitated the development of deep convolutional neural networks~\citep{He:16} as well as Transformers~\citep{Vaswani:17}. Transformers fundamentally drive Large Language Models (LLMs) such as GPT-3~\citep{Brown:20short}, ChatGPT~\citep{Schulman:22short}, GPT-4~\citep{Achiam:23gpt4short}, Megatron-LM~\citep{Shoeybi:19}, Gopher~\citep{Rae:21short}, ERNIE 3.0 Titan~\citep{Wang:21arxivshort}, GLaM~\citep{Du:21glamshort}, Chinese M6~\citep{Lin:21}, AlexaTM 20B~\citep{Soltan:22}, OPT~\citep{Zhang:22opt}, Chinchilla~\citep{Hoffmann:22short}, BLOOM~\citep{Scao:22arxivshort}, GLM-130B~\citep{Zeng:22arxivshort}, LaMDA~\citep{Thoppilan:22short}, PaLM~\citep{Chowdhery:22short}, Llama~\citep{Touvron:23arxiv}, Gemini~\citep{GeminiTeam:23arxiv,Reid:24arxiv}.

\section{Experiments}
\label{sec:Exp}

This section presents an empirical assessment of xLSTM, positioning it alongside established approaches in the realm of language modeling. Section~\ref{sec:ExpSyntethic} is dedicated to probing the distinct capacities of xLSTM through synthetic task evaluation. Section~\ref{sec:ExpComparison} explores the validation set perplexity of multiple language modeling systems, each trained on a corpus comprising 15B tokens sourced from SlimPajama~\citep{Soboleva:23}, and includes ablation analyses of xLSTM. Subsequently, we characterize how model performance scales following methodologies similar to those in \citet{Kaplan:20} and \citet{Brown:20short}. 

Section~\ref{sec:ExpLanguage} features a comprehensive language modeling analysis, where xLSTM is compared against top-performing contenders identified in Section~\ref{sec:ExpComparison}, this time using a substantially larger dataset with 300B tokens from SlimPajama~\citep{Soboleva:23}. Our study proceeds in several stages: we initially explore the generalization of these models to extended context lengths; next, we scrutinize their validation perplexity and effectiveness in various downstream tasks~\citep{Sutawika:24}; thereafter, we measure their performance on 571 text domains within the PALOMA language benchmark~\citep{Magnusson:23arxivshort}; finally, we reexamine scaling patterns, now considering a twentyfold increase in training data.

\label{sec:xlstm_blocks_notation}
Throughout all experimental setups, we adopt the notation xLSTM[$a$:$b$] to denote the proportion $a/b$ between mLSTM-based and sLSTM-based xLSTM blocks. As an illustration, xLSTM[7:1] indicates a configuration where seven of the eight blocks utilize mLSTM and one block utilizes sLSTM. When considering a typical total of 48 blocks, this setup results in 42 blocks based on mLSTM and 6 blocks based on sLSTM. Additionally, all experimental models employ up-projection blocks prior to mLSTM and following sLSTM, respectively.

\subsection{Synthetic Tasks and Long Range Arena}
\label{sec:ExpSyntethic}
We begin by evaluating the performance of xLSTM's novel exponential gating mechanism combined with memory mixing using formal language tasks, as described by~\citep{Deletang:23}. Subsequently, we investigate how xLSTM's recently introduced matrix memory handles the Multi-Query Associative Recall task following~\citep{Arora:23arxiv}. Lastly, we measure xLSTM's capability for long-sequence processing within the Long Range Arena benchmark, referencing~\citep{Tay:21}.

\paragraph{Evaluation of xLSTM's Exponential Gating and Memory Mixing Capabilities.} We evaluate the performance of xLSTM's recently introduced exponential gating mechanism combined with memory mixing, which is theoretically poised to address state tracking tasks~\citep{Merrill:24, Merrill:23}. Formal language benchmarks from \citet{Deletang:23} are modified and expanded to support multi-length training, facilitating extrapolation over sequence lengths. More comprehensive information about the tasks and expanded experimental findings can be found in Appendix~\ref{sec:appExpSynthetic-formalLang}. We benchmark xLSTM against a variety of models, including Transformers, State Space Models, and traditional Recurrent Neural Networks. Model accuracy is calculated based exclusively on task-relevant tokens, with performance normalized on a scale from 0 (indicating random predictions) to 1 (reflecting perfect accuracy). Comparisons are conducted using architectures with two blocks for the following variants: xLSTM[0:1] (solely sLSTM), xLSTM[1:0] (only mLSTM), xLSTM[1:1], Llama, Mamba, RWKV, Retention, Hyena, LSTM, and LSTM embedded within Transformer blocks (LSTM (Block)). Figure~\ref{fig:formal-main} illustrates the outcomes from these experiments. Notably, architectures such as Transformers and State Space Models that lack memory mixing capabilities (thus unable to perform state tracking) are incapable of solving certain tasks—such as those involving regular grammars exemplified by the parity task. This observation aligns with previous reports indicating that Transformers and State Space Models are intrinsically less expressive compared to RNNs~\citep{Merrill:24, Merrill:23, Deletang:23}.

\paragraph{Assessing xLSTM's Memory Capacity via Associative Recall Tasks.} This study investigates the matrix memory mechanism introduced in xLSTM by evaluating its performance on the Multi-Query Associative Recall task~\citep{Arora:23arxiv}, which measures the ability to store and retrieve key-value associations sampled from a sizable vocabulary within sequences. To further challenge the recall mechanism beyond the initial setup, the task is scaled by increasing both the count of key-value pairs to 256 and the context window up to 2048 tokens, thus providing a more rigorous probe of memory limitations across architectures. The experiment includes two-block variants from Llama, Mamba, RWKV-5, RWKV-6, xLSTM[1:1], and xLSTM[1:0]. Performance is quantified by the proportion of successfully retrieved pairs. Transformer-based models such as Llama, characterized by memory scaling exponentially with coding dimensionality~\citep{Ramsauer:21}, serve as the benchmark for this challenge. The empirical outcomes are depicted in Figure~\ref{fig:mqar-main}.

Among non-Transformer architectures, xLSTM[1:1] exhibits the highest accuracy, even outperforming others in the regime of smaller model sizes. Notably, including the sLSTM block not only preserves but also amplifies memory capacity, an effect that becomes especially prominent when handling the largest associative load (256 pairs). Further experimental details and analyses are given in Appendix~\ref{sec:appExpSynthetic-mqar}; extrapolation experiments suggest that xLSTM's memory improvements generalize well to contexts exceeding those present during training.

\paragraph{Test of xLSTM's Long Context Capabilities on Long Range Arena.} In order to evaluate the ability of xLSTM to process lengthy sequences and extended contexts, we benchmark various approaches on the Long Range Arena~\citep{Tay:21}. xLSTM consistently achieves high results across each task, indicating that its architecture is highly effective for addressing diverse long context challenges.

For more details, see Appendix~\ref{sec:appExpSynthetic-lra}.

\subsection{Method Comparison and Ablation Study}
\label{sec:ExpComparison}

In order to explore the central inquiry of this work—namely, the capabilities of our proposed LSTM variants when applied to large-scale language modelling—we conduct experiments training xLSTMs, Transformers, State Space Models, and additional approaches using 15B tokens drawn from SlimPajama under consistent auto-regressive conditions. Model performance is evaluated on the validation dataset, and we further carry out ablation studies focused on xLSTM architectures.

\paragraph{Benchmarking xLSTM Against Alternative Approaches.} To facilitate comparison, each model is trained on a dataset comprising 15B tokens sourced from SlimPajama~\citep{Soboleva:23}. Model performance is assessed through perplexity computed on the validation subset. The comparison includes several approaches: the newly introduced xLSTM, GPT-3 (Transformer)~\citep{Brown:20short}, Llama (Transformer)~\citep{Touvron:23arxiv}, H3 (SSM)~\citep{Fu:23}, Mamba (SSM)~\citep{Gu:24arxiv}, RWKV-4 (RNN)~\citep{Peng:23arxivshort}, RWKV-5 (RNN)~\citep{Peng:24arxivshort}, RWKV-6 (RNN)~\citep{Peng:24arxivshort}, GLA (linear Transformer)~\citep{Yang:23arxiv}, HGRN (RNN)~\citep{Qin:23}, HGRN2 (RNN)~\citep{Qin:24arxiv}, RetNet (linear Transformer)~\citep{Sun:23arxiv}, Hyena (linear Transformer)~\citep{Poli:23}, and the xLSTM variants xLSTM[1:0] and xLSTM[7:1] as detailed in Section~\ref{sec:xlstm_blocks_notation}. 

All models utilize mixed precision training, except for RWKV-5, RWKV-6, GLA, and HGRN2, where mixed precision was implemented differently and did not rely on the PyTorch automatic mixed precision feature (see Appendix Section~\ref{sec:appGenTrainProc} for further information). The reviewed methods are classified into three categories: (a) Transformers, (b) State Space Models (SSMs), and (c) Recurrent Neural Networks (RNNs), along with linear Transformers. Here, linear Transformers are methods that replace standard Transformer attention with a linear mechanism. Model architectures are set to match a 350M parameter GPT-3 in terms of scale, adopting a 1024-dimensional embedding and 24 residual blocks. Notably, GPT-3 is unique in employing weight sharing between token and output embeddings, resulting in a smaller parameter count. 

Findings, presented in Table~\ref{tab:spaj15b_model_results}, reveal that xLSTM achieves superior validation perplexity relative to all other compared models. The Appendix~\ref{sec:appExpComparison} provides additional discussion and methodology. Scaling patterns for these experiments are illustrated in Figure~\ref{fig:temp_sclaw15B}, suggesting that xLSTM continues to excel as model size increases.

\paragraph{Ablation Studies.}
Table~\ref{tab:spaj15b_model_results} along with Figure~\ref{fig:temp_sclaw15B} illustrate xLSTM’s superior performance in language modeling tasks, specifically when trained on 15B tokens drawn from the SlimPajama dataset. In order to analyze the incremental contributions from the modifications leading from LSTM to xLSTM, we progressively transition a standard LSTM architecture toward the full xLSTM design. The initial adjustment involves embedding LSTM layers within a pre-LayerNorm residual framework. Following this, the architecture is modified by appending a post up-projection module. In the final stage, exponential gating mechanisms and matrix memory components are incorporated.

The results are shown in Table~\ref{tab:ablstudies} (top). 

Ablation experiments credit the marked boost in performance to the synergy between exponential gating and the matrix memory component. Furthermore, recognizing the critical role gating mechanisms play in RNNs and State Space Models, we systematically remove or vary different gating types to analyze their effects. As shown in Table~\ref{tab:ablstudies} (bottom), making each gate both learnable and responsive to input progressively enhances model effectiveness. Further detailed investigations regarding specific backbone modules are presented in Appendix~\ref{sec:appAblDetails}.

\subsection{xLSTM as Large Language Model}
\label{sec:ExpLanguage}

In this work, we undertake extensive language modeling experiments to explore xLSTM as a large language model (LLM). To this end, we scale up the training data, utilizing 300B tokens sourced from SlimPajama. This token count is consistent with datasets employed in prior studies, including Mamba~\citep{Gu:24arxiv} and Griffin~\citep{De:24arxiv}.

For comparative analysis, xLSTM is evaluated alongside RWKV-4, Llama, and Mamba—each representing a distinct method class as outlined in Section~\ref{sec:ExpComparison}. RWKV-4 is chosen as the RNN exemplar because suitable training precision configurations for RWKV-5, RWKV-6, and HGRN2 were only established after initiating the 300B token training runs (details can be found in Appendix~\ref{sec:appGenTrainProc}).

We experiment with a range of model sizes (125M, 350M, 760M, 1.3B) and examine their ability to extrapolate to longer sequences. Their validation performance is measured, and further assessment is carried out on downstream tasks. Additionally, their language modeling abilities are benchmarked across 571 text domains in PALOMA, followed by analysis of their scaling law properties.

\paragraph{Sequence Length Extrapolation.}
\label{sec:spaj300B_ctx_extrapolation}
We begin by investigating the sequence length generalization abilities of xLSTM, RWKV-4, Llama, and Mamba, each at the 1.3B parameter scale. Models are trained with sequences of length 2048 and subsequently evaluated on sequences as long as 16384. The outcomes are presented in Figure~\ref{fig:spaj300B_seq_extrapolation}. Notably, xLSTM consistently achieves low perplexity across extended contexts, whereas competing architectures degrade more rapidly.

\paragraph{Validation Perplexity and Downstream Tasks.}
\label{sec:spaj_downstream_eval}
Additionally, for every considered model size, we assess the capabilities of xLSTM, RWKV-4, Llama, and Mamba on next token prediction accuracy using the SlimPajama validation set, as well as their performances on a suite of downstream tasks centered on common sense reasoning.

The third column in Table~\ref{tab:lmeval} presents the validation set perplexities achieved by various approaches. xLSTM[1:0] and xLSTM[7:1] outperform other models across every size with regard to validation perplexity. Performance metrics for downstream tasks are shown in the remaining columns of Table~\ref{tab:lmeval}. Across nearly all tasks and model sizes, xLSTM consistently delivers the strongest results; Mamba surpasses xLSTM only in selected ARC task instances. Further information is provided in Appendix~\ref{sec:appExpLanguage}.

\paragraph{Performance on PALOMA Language Tasks.} We additionally evaluate the efficacy of next token prediction across all model sizes for xLSTM, RWKV-4, Llama, and Mamba on PALOMA language tasks~\citep{Magnusson:23arxivshort}. The evaluation metric employed is perplexity on next token prediction spanning 571 distinct textual domains, encompassing sources from nytimes.com to Reddit’s r/depression.
Table~\ref{tab:ppleval} presents perplexity results categorized into two sections: language modeling benchmarks (first seven columns) and more granular domain benchmarks (last five columns). On these language tasks, xLSTM[1:0] demonstrates superior performance relative to xLSTM[7:1].

Across 568 of the 571 domains (constituting 99.5\%), xLSTM[1:0] achieves lower perplexity scores compared to Mamba; against Llama, it obtains lower perplexity in 486 domains (85.1\%); and versus RWKV-4, it exceeds performance in 570 domains (99.8\%). Detailed results are available in Appendix~\ref{sec:appExpLanguage}.

\paragraph{Scaling Laws.} Fourth, we investigate the power-law scaling behavior, which provides a means to predict model performance at increased scale~\citep{Kaplan:20,Brown:20short}. Figure~\ref{fig:temp_sclaw300B} illustrates this scaling trend. Although all models demonstrate comparable scaling trajectories, they differ in their baseline performance. RWKV-4 yields the lowest results, with Llama and Mamba subsequently improving upon it. xLSTM achieves superior outcomes compared to Mamba, with the difference between xLSTM and Mamba paralleling that between Mamba and Llama. This observed scaling suggests that as model sizes grow, xLSTM maintains a competitive edge over both Transformers and State-Space models.

\textbf{Generation Times and Maximal Throughput.} In Figure~\ref{fig:inference_time_generation}, we present measurements of text generation latency for xLSTM variants at the 1.3B parameter scale. Additionally, the maximal throughput results are displayed in Figure~\ref{fig:inference_time_decoding_throughput} (left). For these benchmarks, we evaluate against similarly-sized implementations of Mamba, Llama, and RWKV from HuggingFace, ensuring that the Llama model employs a static key-value cache. During our experimental runs, we observed that neither full cache compilation for Transformer-based models nor \texttt{torch.compile} support for Mamba functioned reliably. All generation benchmarks were conducted with batch size set to 1 and pre-fill length of 16, a setting that particularly benefits Transformer architectures.

From Figure~\ref{fig:inference_time_generation}, it is evident that xLSTM, along with other recurrent designs such as Mamba and RWKV-4, exhibit linear computational complexity with respect to input, in contrast with the quadratic time scaling observed for Llama. Regarding decoding throughput, we perform evaluations under varying batch sizes and pre-fill conditions specifically for the Llama baseline. As highlighted by Figure~\ref{fig:inference_time_decoding_throughput} (right), the constant memory characteristics of xLSTM enable it to scale efficiently to significantly higher batch sizes compared to Llama, culminating in the highest overall throughput among the models examined.

\section{Limitations}

(i) Unlike mLSTM, the memory mixing approach used by sLSTM limits parallelizable computation, making rapid parallel execution impractical. Despite this, we implemented a custom CUDA kernel for sLSTM, which operates at a speed less than twice that of our mLSTM parallel version. (ii) The mLSTM CUDA kernels lack optimization, resulting in execution roughly four times slower than FlashAttention or the scan mechanism in Mamba. Enhanced CUDA kernels akin to FlashAttention could likely improve performance. (iii) mLSTM’s matrix memory incurs significant computational demands because it requires handling $d \times d$ matrices. However, memory updating and access remain parameter-free and are amenable to standard parallel matrix operations, so the extra elapsed time is negligible. (iv) Gate initialization, particularly for the forget gates, needs to be handled with attention. (v) The matrix memory structure in mLSTM is independent of sequence length, but expanding the context window could overload the memory with longer sequences. Despite this, the approach is feasible for contexts as large as 16k tokens, as discussed in Section~\ref{sec:spaj300B_ctx_extrapolation}. (vi) Given the high computational requirements for large-scale language model experiments, architectural and hyperparameter optimization was not exhaustive, especially for larger xLSTM variants. Comprehensive optimization will be necessary for xLSTM to fully maximize its capabilities.

\section{Conclusion}
We have partially addressed our central inquiry: What are the outcomes of scaling LSTM architectures up to billions of parameters in language modeling? Our findings demonstrate that such scaling enables performance on par with contemporary approaches, including Transformers and State Space Models. By introducing exponential gating with memory mixing and an innovative memory structure, we extend LSTM to the xLSTM framework. In comparative language modeling evaluations, xLSTM displays competitive results against advanced methods such as Transformers and State Space Models. Observed scaling laws suggest that further increasing xLSTM model size could position it as a formidable alternative to Transformer-based Large Language Models. Moreover, xLSTM exhibits promising potential to influence other domains, notably Reinforcement Learning, Time Series Prediction, and the modeling of physical phenomena.

\end{document}